\documentclass[RNAAS]{aastex7}

\begin{document}
\title{SXS: A Reproducible Pipeline for Kepler Transit Recovery and Candidate Vetting}
\author{Rasya Andrean}
\affiliation{Science Experimental Technologies}
\email{scix.official@gmail.com}

\begin{abstract}
SCIX Exoplanet Search (SXS) is a reproducible pipeline combining Box Least
Squares (BLS) transit detection with machine-learning ranking and independent
candidate vetting in public Kepler photometry. In the version 1 benchmark, BLS
recovered 15 of 36 eligible confirmed planets in its five highest peaks, while
the Random Forest (RF) retained 12 of 36 end to end (33.33\%) with a
candidate-level false-positive rate (FPR) of 0.070. The scaled version 2 search
recovered 227 of 434 eligible planets with BLS (52.30\%) and then examined a
deterministic sample of 250 targets without cumulative-KOI or
confirmed-Kepler-name history. Independent review of the resulting shortlist
found no strong candidate, one weak candidate, and 19 likely false positives;
SXS makes no planet-discovery claim.
\end{abstract}

\keywords{Exoplanet detection (489) --- Transit photometry (1709) ---
Astronomy software (1855)}

\section{Data and Methods}
Confirmed-planet parameters came from the NASA Exoplanet Archive
\texttt{pscomppars} table, and official false-positive labels came from the
cumulative KOI table \citep{nasaarchive}. Kepler long-cadence light curves were obtained from MAST
with Lightkurve \citep{lightkurve2018}. After quality filtering, segment
normalization, and Savitzky--Golay detrending, BLS searched periods of
0.5--50 days and retained five distinct peaks per target. A 13-feature RF
ranked peaks using BLS and folded-light-curve diagnostics. Five-fold
\texttt{StratifiedGroupKFold} evaluation grouped all signals by target. The RF
threshold of 0.221107 maximized precision subject to out-of-fold recall of at
least 0.90. Independent validation excluded RF scores from its evidence rules
and used 1,000 segment-wise circular-shuffle BLS searches per target, odd/even
and secondary-eclipse tests following DR25-style diagnostics
\citep{thompson2018}, limb-darkened \texttt{batman} fits
\citep{kreidberg2015}, Gaia DR3 scene checks \citep{gaiaarchive}, TESS
photometry \citep{tessarchive}, and an ExoFOP-derived TOI lookup
\citep{toidocs}.

\section{Results}
The original benchmark contained 36 planets inside the fixed search range.
BLS recovered 15/36 (41.67\%); RF classification retained 12/15 recovered
signals, giving end-to-end recovery of 12/36 (33.33\%). Candidate-level RF
precision was 0.632, recall was 0.800, and FPR was 0.070. In the scaled
benchmark, BLS recovered 227/434 planets (52.30\%). At the exploratory
manual-review threshold, RF v2 achieved precision 0.412, recall 0.903, and FPR
0.146 from target-grouped out-of-fold predictions. The FPR is derived from 292
false passes among 2,000 negative peaks.

\begin{deluxetable}{lcccc}
\tablecaption{Pipeline performance and validation outcome\label{tab:results}}
\tablehead{\colhead{Stage} & \colhead{Precision} & \colhead{Recall} &
\colhead{FPR} & \colhead{Recovery or outcome}}
\startdata
v1 BLS proposal set & 0.130 & 1.000 & 1.000 & 15/36 (41.67\%) \\
v1 RF, threshold 0.5 & 0.632 & 0.800 & 0.070 & 12/36 (33.33\%) \\
v2 BLS proposal set & \nodata & \nodata & \nodata & 227/434 (52.30\%) \\
v2 RF, threshold 0.221107 & 0.412 & 0.903 & 0.146 & 205 TP; 292 FP \\
Independent vetting & \nodata & \nodata & \nodata & 0 strong; 1 weak; 19 likely FP
\enddata
\end{deluxetable}

All 250 screening targets completed acquisition, preprocessing, and blind BLS
search, producing 1,250 peaks. RF v2 flagged 151 peaks; 110 also passed the
available preliminary odd/even, phase-0.5 secondary, and moment-centroid
checks. The 20 highest-ranked signals, spanning 14 KIC targets, were frozen
before independent validation.

No shortlisted signal met the strong empirical-FAP threshold of 0.01. The sole
weak signal, KIC 8300900-r1, has period 5.090289 days and empirical BLS FAP
20/1,001 = 0.01998. It passed the configured odd/even, secondary, morphology,
physical-size, and Gaia checks, but lacked TESS period support and is not a
confirmed exoplanet. Nineteen signals were classified as likely false positives
under prespecified key-failure rules. No position-matched public TOI record was
found for the 14 targets, but catalog absence was assigned no positive
evidential weight.

SXS demonstrates the separation between ranking performance and independent
evidence: an RF can prioritize transit-like morphology while most high-ranked
signals fail a target-specific null test or another independent check. The
result is a reproducible methodology and negative-result case study rather
than a validation catalog.

\section{Limitations and Availability}
The benchmark and bounded 250-target search are too small and selected to
support occurrence-rate inference. The RF threshold is an exploratory review
threshold, not a calibrated planetary probability, and the empirical FAP is
conditional on the shuffle scheme, detrending, sampling, and search grid rather
than a VESPA-style Bayesian false-positive probability. Only four Kepler
products per search target were analyzed, and no pixel-level analysis,
spectroscopy, or new physical follow-up was performed.

Software, configurations, catalog snapshots, metrics, and validation products
are available at \url{https://github.com/Science-Experimental-Technologies/Exoplanet-Search}
(software release v1.2.0). The archived research release is v1.0.0. A permanent
archive DOI will be inserted here only after a Zenodo record is published.

\acknowledgments
This work was independently funded by Rasya Andrean and Urus Foundation. It
made use of public services and data from the NASA Exoplanet Archive, MAST,
Gaia, TESS, and ExoFOP.

\bibliography{references}
\bibliographystyle{aasjournal}
\end{document}
